\chapter{Boltzmann Sampling}
\label{ch13}

\section{Motivation}
\label{ch13:motivation}

Suppose we want to draw a uniformly random element of $\mathcal{A}_n$, the set of all structures in a combinatorial class $\mathcal{A}$ of size exactly $n$. If $\mathcal{A}$ is described by a symbolic specification, the classical approach is the \emph{recursive method} of Nijenhuis and Wilf \cite{NijenhuisWilf1978}, later systematized by Flajolet, Zimmermann and Van Cutsem \cite{FlajoletZimmermannVanCutsem1994}. The idea is to precompute a table of values $[z^k]A(z)$ for $k \leq n$, then walk down the recursive structure of the specification, at each node making a random choice with probabilities proportional to those table entries. This is exact and elegant, but it has costs: storing the table requires $O(n^2)$ space (one entry per size per subclass), and generating a single sample takes $O(n \log n)$ arithmetic operations once the table is built.

For large $n$, or for classes whose specifications are complex, this becomes unwieldy. Duchon, Flajolet, Louchard and Schaeffer \cite{DuchonFlajoletLouchardSchaeffer2004} proposed an alternative: rather than insisting on \emph{exact} size $n$, allow the size to be random, governed by a parameter $x$. By tuning $x$ one can concentrate the size distribution near a target $n$, and by combining this with rejection sampling one recovers a uniform sampler for exact size $n$ at a dramatically reduced cost — $\Theta(n)$ expected time for well-behaved classes, with no precomputed tables.

The approach is called \emph{Boltzmann sampling}, and as we shall see, the name is not accidental. The probability weights assigned to structures are precisely Boltzmann weights from statistical physics. This identification — generating function equals partition function — is the central observation of this chapter, and it will be the bridge to Chapter 14, where LLM temperature decoding is read in exactly the same terms.

\section{Boltzmann Samplers}
\label{ch13:boltzmann-samplers}

\subsection{Definition}
\label{ch13:def}

Let $\mathcal{A}$ be a combinatorial class with generating function $A(x) = \sum_{n \geq 0} a_n x^n$, where $a_n = |\mathcal{A}_n|$, and let $\rho_A$ denote the radius of convergence of $A$.

\begin{definition}
\label{ch13:def-sampler}
Fix $x \in (0, \rho_A)$. A \emph{(free) Boltzmann sampler} $\Gamma A(x)$ is a random procedure that produces an object $\gamma \in \mathcal{A}$ with probability
\[
    \Pr_x(\gamma) = \frac{x^{|\gamma|}}{A(x)}.
\]
\end{definition}

Several observations are immediate. First, this is a valid probability distribution: summing over all $\gamma \in \mathcal{A}$,
\[
    \sum_{\gamma \in \mathcal{A}} \frac{x^{|\gamma|}}{A(x)}
    = \frac{1}{A(x)} \sum_{n \geq 0} a_n x^n = 1.
\]
Second, all objects of the \emph{same size} $n$ receive equal probability $x^n / A(x)$. So conditioning on size, $\Gamma A(x)$ yields a uniform element of $\mathcal{A}_n$. This is the key uniformity property that makes Boltzmann samplers useful for random generation.

Third, the size $N = |\gamma|$ is a random variable. Its distribution is
\[
    \Pr_x(N = n) = \frac{a_n x^n}{A(x)},
\]
which is the probability generating function (PGF) model with weight $a_n x^n$ normalized by $A(x)$. The expected size is
\[
    \mathbb{E}_x[N] = \sum_{n \geq 0} n \cdot \frac{a_n x^n}{A(x)} = \frac{x A'(x)}{A(x)}.
\]

This last formula is important: as $x$ increases toward $\rho_A$, the expected size grows (often without bound), so we can tune the expected size of our random sample by adjusting $x$.

\subsection{Construction from Symbolic Specifications}
\label{ch13:construction}

The power of the Boltzmann framework lies in how naturally it interacts with the symbolic method. If $\mathcal{A}$ is defined by a combinatorial specification, the corresponding Boltzmann sampler is built recursively from the same structure.

We record the four main constructors; proofs are straightforward from the definition.

\paragraph{Disjoint union.} Suppose $\mathcal{A} = \mathcal{B} + \mathcal{C}$, so $A(x) = B(x) + C(x)$. Then $\Gamma A(x)$ proceeds as follows: flip a biased coin that lands on $\mathcal{B}$ with probability $B(x)/A(x)$ and on $\mathcal{C}$ with probability $C(x)/A(x)$; if $\mathcal{B}$ is chosen, return $\Gamma B(x)$; otherwise return $\Gamma C(x)$.

To verify: an object $\gamma \in \mathcal{B}$ is produced with probability
\[
    \frac{B(x)}{A(x)} \cdot \frac{x^{|\gamma|}}{B(x)} = \frac{x^{|\gamma|}}{A(x)},
\]
and similarly for $\mathcal{C}$, so the definition is satisfied.

\paragraph{Cartesian product.} Suppose $\mathcal{A} = \mathcal{B} \times \mathcal{C}$, so $A(x) = B(x) \cdot C(x)$. Then $\Gamma A(x)$ independently draws $\beta \sim \Gamma B(x)$ and $\gamma \sim \Gamma C(x)$, and returns the pair $(\beta, \gamma)$.

Indeed, $\Pr_x((\beta, \gamma)) = \frac{x^{|\beta|}}{B(x)} \cdot \frac{x^{|\gamma|}}{C(x)} = \frac{x^{|\beta|+|\gamma|}}{A(x)}$, as required.

\paragraph{Sequence.} Suppose $\mathcal{A} = \SEQ(\mathcal{B})$, so $A(x) = 1/(1 - B(x))$, valid when $B(x) < 1$. Draw a length $K$ from the geometric distribution $\Pr(K = k) = (1 - B(x)) B(x)^k$ for $k \geq 0$, then independently draw $K$ objects from $\Gamma B(x)$, and return the sequence $(\beta_1, \ldots, \beta_K)$.

The probability of producing a specific sequence $(\beta_1, \ldots, \beta_k)$ of total size $n = |\beta_1| + \cdots + |\beta_k|$ is
\[
    (1 - B(x)) B(x)^k \cdot \prod_{i=1}^k \frac{x^{|\beta_i|}}{B(x)}
    = (1 - B(x)) \cdot \frac{x^n}{B(x)^k} \cdot B(x)^k
    = (1 - B(x)) x^n = \frac{x^n}{A(x)},
\]
confirming the Boltzmann property.

\paragraph{Set and cycle.} For $\mathcal{A} = \SET(\mathcal{B})$ and $\mathcal{A} = \CYC(\mathcal{B})$, the constructions involve Poisson distributions and logarithmic-series distributions respectively. Specifically, the number of copies of a given unlabeled component of size $k$ in a random set structure follows a Poisson distribution with mean $(x^k/k!)$, and analogously for cycles. We do not spell out the full surgery here; the reader should consult \cite{FlajoletSedgewick2009}, Chapter VII, for the complete treatment.

\subsection{Tuning the Size and Rejection Sampling}
\label{ch13:rejection}

For a fixed target size $n$, one uses $\Gamma A(x)$ as a subroutine in a rejection scheme: run $\Gamma A(x)$ repeatedly and keep only outputs of size $n$. Since all such outputs are equally likely (by the uniformity property), each kept sample is uniform over $\mathcal{A}_n$.

The question is efficiency. The probability of getting an object of size exactly $n$ is $a_n x^n / A(x)$. To maximize this, one should choose $x$ to maximize $a_n x^n / A(x)$ over $x \in (0, \rho)$, which by a saddle-point argument gives $x^* = x^*(n)$ satisfying $x^* A'(x^*)/A(x^*) = n$, i.e., the expected size equals $n$.

\begin{proposition}
\label{ch13:rejection-complexity}
For combinatorial classes satisfying the smoothness conditions of \cite{DuchonFlajoletLouchardSchaeffer2004} (roughly: $a_n \sim C \rho^{-n} n^{-\alpha}$ for some $\alpha > 0$), rejection sampling with the optimal parameter $x^*(n)$ produces a uniform element of $\mathcal{A}_n$ in $\Theta(\sqrt{n})$ expected trials if $\alpha = 3/2$, and $\Theta(1)$ if $\alpha > 1$. In either case the total expected work is $\Theta(n)$.
\end{proposition}

This should be contrasted with the recursive method's $O(n^2)$ preprocessing: for large classes and large $n$, the Boltzmann approach wins decisively.

\section{The Gibbs Measure Interpretation}
\label{ch13:gibbs}

We now come to the central observation of this chapter.

\subsection{Boltzmann Weights as Gibbs Weights}
\label{ch13:gibbs-weights}

Make the substitution $x = e^{-\beta}$, where $\beta \in \mathbb{R}$ is to be thought of as an \emph{inverse temperature}. The Boltzmann sampler probability becomes
\[
    \Pr_\beta(\gamma) = \frac{e^{-\beta |\gamma|}}{A(e^{-\beta})}.
\]
Define the \emph{energy} of a structure $\gamma$ to be its size: $E(\gamma) = |\gamma|$. Then
\[
    \Pr_\beta(\gamma) = \frac{e^{-\beta E(\gamma)}}{Z(\beta)}, \qquad Z(\beta) = A(e^{-\beta}).
\]

This is exactly the Gibbs (Boltzmann) distribution from statistical mechanics. The normalization $Z(\beta)$ is the \emph{partition function}, and here it is simply the generating function of the class, evaluated at $e^{-\beta}$.

\begin{remark}
\label{ch13:gibbs-correspondence}
\textbf{The generating function is the partition function.} For any combinatorial class $\mathcal{A}$ with generating function $A(z)$ and energy function $E(\gamma) = |\gamma|$, the Gibbs distribution at inverse temperature $\beta$ is
\[
    \Pr_\beta(\gamma) = \frac{e^{-\beta |\gamma|}}{A(e^{-\beta})},
\]
and the partition function is $Z(\beta) = A(e^{-\beta})$. This identification is not a formal analogy — it is an equality of formulas. The analytic-combinatorics side (generating functions, singularities, coefficients) and the statistical-physics side (partition functions, free energy, phase transitions) are the same mathematical object viewed through different lenses. This is the single most important observation of this chapter.
\end{remark}

\subsection{The Critical Temperature}
\label{ch13:critical-temp}

The radius of convergence $\rho_A$ of $A(z)$ corresponds to a \emph{critical inverse temperature}
\[
    \beta_c = -\log \rho_A.
\]
For $\beta > \beta_c$ (i.e., $e^{-\beta} < \rho_A$), the partition function $Z(\beta) = A(e^{-\beta})$ is finite, and the Gibbs measure is a well-defined probability distribution on $\bigcup_{n \geq 0} \mathcal{A}_n$.

For $\beta < \beta_c$ (i.e., $e^{-\beta} > \rho_A$), the series $A(e^{-\beta})$ diverges, so no normalizable Gibbs measure exists on the whole class. One can still define a uniform distribution on each fixed $\mathcal{A}_n$ separately (the recursive method does exactly this), but there is no consistent extension to a single measure across all sizes.

At $\beta = \beta_c$, whether the partition function converges depends on the growth rate of $a_n$. For many natural classes, $a_n \sim C \rho^{-n} n^{-\alpha}$ with $\alpha = 3/2$ (trees, planar maps, etc.), in which case $Z(\beta_c) = A(\rho)$ diverges since $\sum n^{-3/2} < \infty$ but barely so. The singular sampler at $x = \rho$ is a borderline case that produces a well-defined but heavy-tailed size distribution, which feeds into the rejection scheme of Proposition~\ref{ch13:rejection-complexity}.

\section{Free Energy and the Thermodynamic Limit}
\label{ch13:free-energy}

For finite-system statistical mechanics, it is standard to study the partition function restricted to systems of size $n$:
\[
    Z_n(\beta) = a_n \, e^{-\beta n} = [z^n] A(z) \cdot e^{-\beta n}.
\]
This is simply $a_n$ rescaled, and it is the unnormalized probability of drawing a structure of size $n$ from the Gibbs distribution.

The \emph{free energy per unit size} is
\[
    f(\beta) = -\lim_{n \to \infty} \frac{1}{n} \log Z_n(\beta)
    = \beta - \lim_{n \to \infty} \frac{1}{n} \log a_n.
\]
By the standard theory of exponential growth rates, if $a_n \sim C \rho^{-n} n^{-\alpha}$, then $\lim_{n \to \infty} n^{-1} \log a_n = \log(1/\rho)$. Hence
\[
    f(\beta) = \beta - \log(1/\rho) = \beta + \log \rho.
\]
At the critical temperature $\beta_c = -\log \rho$, the free energy vanishes: $f(\beta_c) = 0$.

This connects directly to the entropy rate of Chapter 9. There we saw that the exponential growth rate $\log(1/R)$ of the coefficient sequence of a generating function is the entropy rate of the associated random process (with $R$ the radius of convergence). Here the same quantity $\log(1/\rho)$ appears as the thermodynamic driving force: the free energy is zero when the inverse temperature is tuned to the entropy rate.

At high temperature ($\beta$ small, $x = e^{-\beta}$ close to $1$), the Gibbs measure strongly favors large structures; the expected size $xA'(x)/A(x)$ diverges as $x \to \rho^-$. At low temperature ($\beta$ large), the measure concentrates on small structures, with $\Pr_\beta(N = 0) \to 1$ as $\beta \to \infty$ if $a_0 > 0$.

\section{Worked Example: Binary Strings}
\label{ch13:example}

\begin{example}
\label{ch13:ex-binary}
Let $\mathcal{A} = \SEQ(\mathcal{Z}_0 + \mathcal{Z}_1)$ be the class of binary strings, where $\mathcal{Z}_0$ and $\mathcal{Z}_1$ are atoms of size 1. The generating function is
\[
    A(z) = \frac{1}{1 - 2z}, \qquad \rho_A = \frac{1}{2}.
\]

\paragraph{The sampler.} To run $\Gamma A(x)$, we use the sequence recipe: draw $K \sim \mathrm{Geom}(1 - 2x)$ (so $\Pr(K = k) = (1-2x)(2x)^k$), then draw each of the $K$ bits independently and uniformly from $\{0, 1\}$. The output is a uniformly random binary string of length $K$.

\paragraph{Size distribution.} The length $K$ has distribution
\[
    \Pr_x(K = n) = \frac{2^n x^n}{A(x)} = (1 - 2x)(2x)^n,
\]
which is geometric with success parameter $1 - 2x$. The expected length is
\[
    \mathbb{E}_x[K] = \frac{x A'(x)}{A(x)} = \frac{x \cdot 2/(1-2x)^2}{1/(1-2x)} = \frac{2x}{1 - 2x}.
\]
Setting this equal to a target $n$ and solving gives $x^*(n) = n/(2(n+1))$, which approaches $\rho = 1/2$ as $n \to \infty$.

\paragraph{Partition function.} Substituting $x = e^{-\beta}$:
\[
    Z(\beta) = A(e^{-\beta}) = \frac{1}{1 - 2e^{-\beta}}.
\]
This is finite for $\beta > \beta_c = \log 2$. The partition function restricted to size $n$ is $Z_n(\beta) = 2^n e^{-\beta n}$, and the free energy per bit is $f(\beta) = \beta - \log 2$.

\paragraph{Uniformity check.} Among all binary strings of length $n$, there are $2^n$ of them, each with probability $x^n / A(x)$, so each has equal probability $x^n (1-2x) / (2x)^n = (1-2x)/2^n \cdot (1/x)^n \cdot x^n$. Wait — more directly, among strings of length $n$, $\Pr_x(\text{specific string}) = x^n / A(x)$, and since all $2^n$ such strings have the same weight, they are equiprobable, as required.
\end{example}

\begin{example}
\label{ch13:ex-catalan}
For comparison, consider plane binary trees (Catalan structures), with generating function
\[
    C(z) = \frac{1 - \sqrt{1 - 4z}}{2z}, \qquad \rho_C = \frac{1}{4},
\]
so $[z^n] C(z) = \binom{2n}{n}/(n+1) \sim 4^n / (n^{3/2} \sqrt{\pi})$. The partition function at inverse temperature $\beta$ is $Z(\beta) = C(e^{-\beta})$, finite for $\beta > \beta_c = \log 4$.

The free energy per node is $f(\beta) = \beta - \log 4$. From the statistical-physics side, the $n^{-3/2}$ factor in the coefficient asymptotics appears as a sub-exponential correction to $Z_n(\beta)$: $Z_n(\beta_c) = [z^n]C(z) \sim C n^{-3/2}$. This is precisely the $n^{-3/2}$ singularity exponent that Chapter 5 derived via singularity analysis of $C(z)$ at $z = 1/4$. The same fact — viewed there as a transfer lemma for square-root singularities, viewed here as a sub-leading correction to the free energy — is the same computation.
\end{example}

\section{Preview of Chapter 14}
\label{ch13:preview}

The machinery developed here admits a direct generalization that connects to language model temperature decoding.

Replace the combinatorial class $\mathcal{A}$ with the set $\mathcal{W}^*$ of finite strings over a vocabulary $\mathcal{V}$, and replace the uniform-within-size distribution of the Boltzmann sampler with the distribution $\mu$ defined by an autoregressive LLM (as in Chapter 10). Define the energy of a string $w$ by
\[
    E(w) = -\log \mu(w),
\]
the negative log-probability under the model. Then temperature-$T$ sampling produces $w$ with probability proportional to
\[
    \mu(w)^{1/T} = e^{-E(w)/T},
\]
which is a Gibbs distribution at inverse temperature $\beta = 1/T$ with energy $E(w)$.

The partition function is $Z(\beta) = \sum_w \mu(w)^\beta$, a Rényi-type sum over all strings. At $\beta = 1$ (temperature 1), this is just $\sum_w \mu(w) = 1$, always finite. At $\beta > 1$ (low temperature), the measure sharpens; at $\beta < 1$ (high temperature), it flattens toward uniform. At $\beta \to 0$ (infinite temperature), all strings of a given length become equally likely, which is the analogue of the Boltzmann sampler at $x \to 1$.

The analogy is exact: the LLM is a partition function machine. Chapter 14 will develop this correspondence carefully, using the WFA approximation of Chapter 12 to connect the LLM partition function to a rational generating function, and thereby apply the singularity analysis of Chapters 3--5 to the thermodynamics of language model decoding.

\section{Notes and Further Reading}
\label{ch13:notes}

The Boltzmann sampler framework originates in \cite{DuchonFlajoletLouchardSchaeffer2004}, which gives full proofs of all the constructions in Section~\ref{ch13:construction} and the complexity results of Proposition~\ref{ch13:rejection-complexity}. The connection between generating functions and partition functions is classical in combinatorics (it is implicit in Polya's enumeration theorem and the transfer matrix method), but the Boltzmann sampling paper makes it algorithmic.

The recursive method, against which Boltzmann sampling is compared in Section~\ref{ch13:motivation}, is developed in \cite{NijenhuisWilf1978} and \cite{FlajoletZimmermannVanCutsem1994}. For sets and cycles (sketched in Section~\ref{ch13:construction}), the detailed Boltzmann constructions appear in \cite{FlajoletSedgewick2009}, Chapter VII.

The identification of the radius of convergence with the critical inverse temperature, and the free energy computation of Section~\ref{ch13:free-energy}, are standard in the statistical mechanics of polymers and random combinatorial structures; see the monograph of Velenik for a textbook treatment in the physics tradition.

\section{Exercises}
\label{ch13:exercises}

\begin{exercise}
\label{ch13:ex-plane-trees}
Let $\mathcal{T}$ be the class of rooted plane trees, satisfying the specification $\mathcal{T} = \mathcal{Z} \times \SEQ(\mathcal{T})$, so $T(z) = z/(1 - T(z))$, giving $T(z) = (1 - \sqrt{1 - 4z})/2$ and $\rho_T = 1/4$.
\begin{enumerate}
    \item[(a)] Describe the Boltzmann sampler $\Gamma T(x)$ as an explicit random procedure using the recursive specification.
    \item[(b)] Compute $\mathbb{E}_x[N]$ and find the value $x^*(n)$ that makes the expected tree size equal to $n$.
    \item[(c)] Compute the critical inverse temperature $\beta_c$ and the free energy $f(\beta)$ for $\beta > \beta_c$.
    \item[(d)] Explain why the singularity exponent $-3/2$ in $[z^n]T(z) \sim C \cdot 4^n n^{-3/2}$ implies that the Gibbs measure at $\beta = \beta_c$ assigns probability $\Pr_{\beta_c}(N = n) \sim C' n^{-3/2}$ to size-$n$ trees, and find the normalizing constant (up to a multiplicative constant).
\end{enumerate}
\end{exercise}

\begin{exercise}
\label{ch13:ex-temperature}
Let $\mathcal{A}$ be a combinatorial class with generating function $A(z) = \sum_n a_n z^n$ and radius of convergence $\rho$.
\begin{enumerate}
    \item[(a)] Show that the variance of the size under $\Gamma A(x)$ is $\mathrm{Var}_x(N) = x \frac{d}{dx}\left(\frac{x A'(x)}{A(x)}\right)$.
    \item[(b)] For $A(z) = e^z$ (the class of all sets of labeled atoms), compute $\mathbb{E}_x[N]$ and $\mathrm{Var}_x(N)$. What distribution does $N$ follow?
    \item[(c)] In the Gibbs interpretation with $x = e^{-\beta}$, express the variance $\mathrm{Var}_\beta(E)$ of the energy (here $E(\gamma) = |\gamma|$) in terms of $\beta$ and $A$. This quantity is the \emph{heat capacity} $C_V(\beta) = -\partial^2 \log Z / \partial \beta^2$. Verify your formula against the result of part (b).
\end{enumerate}
\end{exercise}
